\chapter{DISCUSIÓN}

Con el experimento ya ejecutado sobre la biblioteca de referencia real, este capítulo revisa qué funcionó del diseño, qué limita los resultados obtenidos y en qué se apartó el proyecto del planteamiento inicial.

\section{Adecuación de la metodología al problema}

Tras el pivote (ADR-004), el diseño experimental se ha mostrado adecuado al objetivo. Concentrar la contribución en una comparación empírica y reproducible de siete métricas de similitud, en lugar de en el desarrollo o el ajuste de un clasificador, responde a la pregunta de investigación sin depender de la disponibilidad ni de la idoneidad de modelos preentrenados para el dominio de los microplásticos.

El \textit{pipeline} de preprocesado reproduce el protocolo de referencia del director con una diferencia máxima absoluta de~0. Es la condición de partida del trabajo: sin ella, cualquier diferencia observada entre métricas podría venir de un error de implementación y no del fenómeno medido.

\section{Cambios respecto al planteamiento inicial}

Respecto al anteproyecto original, el cambio de mayor calado fue el pivote de diseño del ADR-004. La arquitectura prevista al principio, con ingesta mediante Apache Kafka y clasificación con los modelos preentrenados de Open Specy sobre un flujo Spark Structured Streaming, se sustituyó por el diseño actual, centrado en la robustez de las métricas de similitud frente a la degradación de la calidad del espectro. El director y yo acordamos el cambio por tres razones:

\begin{itemize}
    \item La literatura no recoge una comparación sistemática y a escala de las métricas de similitud habituales en \textit{matching} espectral frente a distintos niveles de degradación de la señal. La pregunta nueva estaba sin responder; la anterior, no.
    \item Reutilizar los clasificadores de Open Specy como caja negra habría acotado la contribución del TFM a una capa de ingesta y orquestación, sin evidencia empírica propia sobre el problema de fondo: la fiabilidad del \textit{matching} cuando la señal es de baja calidad.
    \item El diseño nuevo encaja con los datos reales disponibles (espectros Raman de cinco polímeros con distintos tiempos de integración) y con el protocolo de referencia ya validado por el director. Depender de una integración en tiempo real con los instrumentos habría añadido un riesgo de ejecución que este camino evita.
\end{itemize}

El pivote atraviesa toda la memoria: los capítulos de Objetivos, Antecedentes/Estado del arte y Desarrollo están redactados conforme al diseño vigente.

\section{Limitaciones identificadas}

El trabajo tiene cuatro limitaciones:

\begin{itemize}
    \item \textbf{Cobertura incompleta de la base de datos de referencia.} Dos de los cinco polímeros medidos, BioPBS y ECOVIO, no figuran en el catálogo de la \emph{Hagelskjær Microplastic Solution Library~1.1} (28~polímeros y 2~minerales), de modo que su exactitud de identificación es del 0\% con las siete métricas y en todos los niveles de calidad. La causa está en la biblioteca, no en el \textit{pipeline}. El análisis estadístico se restringe por ello a los tres polímeros con verdad-terreno disponible: PET, PLA y PP.
    \item \textbf{Tamaño muestral reducido para el análisis estadístico.} Con solo esos tres polímeros, el test de Friedman dispone de pocos bloques: 39 en acumulación y 24 en degradación. La potencia estadística queda limitada y ninguna diferencia entre métricas sobrevive la corrección de Holm por pares, pese a que la exactitud sí muestra una tendencia.
    \item \textbf{Escalado distribuido validado en local, no en clúster real.} La paralelización con Apache Spark se valida en un entorno de un único nodo. El despliegue en infraestructura distribuida física queda fuera del alcance de este TFM y se plantea como trabajo futuro.
    \item \textbf{Cinco polímeros y una única base de datos de referencia.} Las conclusiones sobre robustez se han obtenido sobre un conjunto acotado de polímeros (PET, PP, PLA, BioPBS y ECOVIO) y una sola biblioteca. Generalizarlas a otras familias de polímeros o a espectroscopía FTIR exigiría repetir el análisis, algo que el \textit{pipeline} admite sin rediseño.
\end{itemize}

\section{Interpretación de los primeros resultados}

Entre los dos mecanismos de degradación el patrón no coincide del todo (Capítulo~4). Con acumulación real la métrica más robusta es Pearson; con degradación paramétrica lo son coseno, euclídea, HQI y SAM, que empatan entre sí. Los dos mecanismos sí coinciden en las dos últimas posiciones: Manhattan y Spearman.

Existe una explicación plausible para esta discrepancia. La acumulación real mejora la relación señal-ruido, pero puede dejar restos de deriva de línea base incluso después de \texttt{airPLS}. Pearson es inmune a esos desplazamientos aditivos; coseno y SAM son invariantes a escala, no a offset. Con la degradación paramétrica esa ventaja se diluye, porque allí el ruido gaussiano se añade sobre un espectro ya limpio y no queda offset que corregir.

La discrepancia tiene una consecuencia práctica: la métrica más adecuada depende del mecanismo de pérdida de calidad, y el resultado no debería trasladarse sin más a otros contextos de ruido instrumental.

\section{Impacto esperado}

La contribución principal es un \textit{pipeline} completo de preprocesado, comparación de métricas y análisis estadístico, validado numéricamente, cubierto por pruebas automatizadas y ya ejecutado sobre la biblioteca de referencia real. De él sale una guía empírica de qué métrica aguanta mejor según el mecanismo de degradación, aplicable tanto al análisis de microplásticos como a otros problemas de \textit{matching} espectral con señal degradada. Conviene mantener dos cautelas: el número de polímeros con verdad-terreno disponible es reducido, y la biblioteca de referencia no cubre todos los materiales medidos.
